% ============================================================
% Notes: Fractional calculus via discrete calculus
% Topics: Grünwald--Letnikov derivative, discrete origin,
%         nonlocality (memory), numerical interpretation
% Source: Herrmann, Fractional Calculus: An Introduction for Physicists
% ============================================================

\section{Fractional calculus using discrete calculus}

\subsection{Discrete derivative and finite differences}
Start from the forward difference operator:
\begin{equation}
(\Delta_h f)(t) := f(t+h)-f(t),
\end{equation}
and the backward difference:
\begin{equation}
(\nabla_h f)(t) := f(t)-f(t-h).
\end{equation}

For integer order $n\in\mathbb{N}$, the $n$th backward difference is
\begin{equation}
\nabla_h^n f(t)
=
\sum_{k=0}^{n}(-1)^k\binom{n}{k}f(t-kh).
\end{equation}

The classical derivative is recovered as
\begin{equation}
f'(t)=\lim_{h\to 0^+}\frac{\nabla_h f(t)}{h}.
\end{equation}

% ------------------------------------------------------------

\section{Grünwald--Letnikov (GL) fractional derivative}

\subsection{Definition}
The Grünwald--Letnikov fractional derivative extends the finite-difference form
from integer $n$ to real $\alpha$ by replacing $\binom{n}{k}$ with $\binom{\alpha}{k}$:
\begin{equation}
\boxed{
(D_{GL}^\alpha f)(t)
=
\lim_{h\to 0^+}\frac{1}{h^\alpha}
\sum_{k=0}^{\lfloor t/h\rfloor}(-1)^k\binom{\alpha}{k}\,f(t-kh).
}
\end{equation}

Here
\begin{equation}
\boxed{
\binom{\alpha}{k}
=
\frac{\Gamma(\alpha+1)}{\Gamma(k+1)\Gamma(\alpha-k+1)}.
}
\end{equation}

\subsection{Consistency with integer derivatives}
If $\alpha=n\in\mathbb{N}$, then
\[
\binom{n}{k}=0\quad \text{for }k>n,
\]
so the sum truncates and GL reduces to the standard $n$th derivative limit.

\subsection{Alternative coefficient form}
Define coefficients
\begin{equation}
c_k^{(\alpha)} := (-1)^k\binom{\alpha}{k}.
\end{equation}
Then
\begin{equation}
(D_{GL}^\alpha f)(t)
=
\lim_{h\to 0^+}\frac{1}{h^\alpha}
\sum_{k=0}^{\lfloor t/h\rfloor}c_k^{(\alpha)}f(t-kh).
\end{equation}

\subsection{Recursive computation of coefficients (useful numerically)}
The binomial coefficients satisfy a recursion:
\begin{equation}
c_0^{(\alpha)}=1,
\qquad
c_k^{(\alpha)}=\left(1-\frac{\alpha+1}{k}\right)c_{k-1}^{(\alpha)}
\quad (k\ge 1).
\end{equation}
This allows efficient evaluation of the GL derivative on a grid.

% ------------------------------------------------------------

\section{Relation to non-locality (memory)}

\subsection{GL derivative is explicitly nonlocal}
The GL derivative at time $t$ uses \emph{all past values}
\[
f(t),\ f(t-h),\ f(t-2h),\dots,f(t-\lfloor t/h\rfloor h).
\]
So the derivative at $t$ depends on the entire history on $[0,t]$.

This is the discrete analogue of Volterra-type memory integrals emphasized in nonlocality theory. :contentReference[oaicite:1]{index=1}

\subsection{Memory kernel behavior (long-range influence)}
For non-integer $\alpha$, the coefficients $c_k^{(\alpha)}$ do not truncate.
They decay slowly with $k$ (power-law type), so distant past contributes.

Hence fractional differentiation encodes \textbf{long memory}:
\begin{equation}
(D^\alpha f)(t)\ \text{depends on}\ \{f(\tau):0\le \tau\le t\}.
\end{equation}

\subsection{Local limit}
In the limit $\alpha\to 1$, the GL derivative approaches the ordinary derivative
(which is local in the continuum limit):
\begin{equation}
\lim_{\alpha\to 1}D_{GL}^\alpha f(t) = f'(t).
\end{equation}

% ------------------------------------------------------------

\section{GL derivative vs Riemann--Liouville / Caputo}

\subsection{Equivalence (conceptual)}
Under suitable smoothness assumptions, the GL derivative coincides with the
Riemann--Liouville derivative:
\begin{equation}
D_{GL}^\alpha f(t) = D_{RL}^\alpha f(t),
\end{equation}
so GL provides a \textbf{discrete definition} of the RL derivative.

\subsection{Why GL is important in applications}
\begin{itemize}
\item It gives an immediate \textbf{finite-difference discretization} for fractional PDEs/ODEs.
\item It makes nonlocality obvious: the derivative is a weighted sum of the entire past.
\item It naturally leads to numerical ``history'' terms (memory cost grows with time).
\end{itemize}

% ------------------------------------------------------------

\section{Practical discretization formula}

On a grid $t_n = nh$, define $f_n=f(t_n)$.
A common GL approximation is:
\begin{equation}
(D^\alpha f)(t_n)
\approx
\frac{1}{h^\alpha}\sum_{k=0}^{n}c_k^{(\alpha)}f_{n-k},
\qquad
c_k^{(\alpha)}=(-1)^k\binom{\alpha}{k}.
\end{equation}

\paragraph{Computational cost.}
This is $O(n)$ work per time step if done directly, since all history values appear.
This is the numerical manifestation of \textbf{nonlocal memory}.

% ------------------------------------------------------------

\section{High-yield summary}

\begin{itemize}
\item Finite differences define integer derivatives; GL extends them to fractional order.
\item Grünwald--Letnikov derivative:
\[
D_{GL}^\alpha f(t)
=
\lim_{h\to 0^+}\frac{1}{h^\alpha}
\sum_{k=0}^{\lfloor t/h\rfloor}(-1)^k\binom{\alpha}{k}f(t-kh).
\]
\item Coefficients use Gamma functions:
\[
\binom{\alpha}{k}=\frac{\Gamma(\alpha+1)}{\Gamma(k+1)\Gamma(\alpha-k+1)}.
\]
\item GL derivative is \textbf{explicitly nonlocal} (depends on entire history).
\item Numerically:
\[
D^\alpha f(t_n)\approx h^{-\alpha}\sum_{k=0}^n c_k^{(\alpha)}f_{n-k},
\]
so fractional dynamics naturally carries memory cost.
\end{itemize}
